%%------------------------------------------------------
% UNIVERSIDADE FEDERAL DE SANTA CATARINA - UFSC
%
% Projeto: Master Mind - Jogo de Lógica e Dedução
% Disciplina: Programação para Web
% Professor: Wyllian B. da Silva
%
% Template: estilo IEEEtran [paper com duas colunas]
%%------------------------------------------------------

\documentclass[journal]{IEEEtran}

%%------------------------------------------------------
%% Packages
%%------------------------------------------------------
\usepackage[T1]{fontenc}
\usepackage[utf8]{inputenc}
\usepackage[dvips]{graphicx}
\usepackage[english,brazil]{babel}
\usepackage{pgf}
\usepackage{epsfig}
\usepackage{graphics}
\usepackage{graphicx}
\usepackage{epstopdf}
\usepackage{float}
\usepackage{eqparbox}
\usepackage{hyphenat}
\usepackage{hyperref}

%%------------------------------------------------------
%% Configurações
%%------------------------------------------------------

%%----------------- Título
\title{Master Mind: Implementação de Jogo Clássico Utilizando React e Next.js com Estética Retrô 8-bit}

\newcommand{\emailautor}{marco.arruda@grad.ufsc.br}
\newcommand{\siglaRevista}{UFSC}
\newcommand{\Revista}{Universidade Federal de Santa Catarina (UFSC)}

%%----------------- Autores
\newcommand{\prenomePrincipal}{Marco}
\newcommand{\nomedomeioPrincipal}{Antônio Machado de}
\newcommand{\sobrenomePrincipal}{Arruda}

\expandafter\newcommand\csname prenome2\endcsname{Caio}
\expandafter\newcommand\csname nomedomeio2\endcsname{César Rodrigues de}
\expandafter\newcommand\csname sobrenome2\endcsname{Aquino}

%%----------------- Autor(es)
\author{\IEEEauthorblockN{\prenomePrincipal~\nomedomeioPrincipal~\sobrenomePrincipal\IEEEauthorrefmark{1},
\prenome{2}~\nomedomeio{2}~\sobrenome{2}\IEEEauthorrefmark{1}}

\IEEEauthorblockA{\IEEEauthorrefmark{1}Universidade Federal de Santa Catarina (UFSC)}

\thanks{\Revista. Correspondência aos autores: \prenomePrincipal~\nomedomeioPrincipal~\sobrenomePrincipal~(email: \emailautor).}}

%%------------------------------------------------------
%% Resumo e Palavras-chave
%%------------------------------------------------------
\IEEEtitleabstractindextext{%
\begin{abstract}
Este trabalho apresenta a implementação de uma versão digital do jogo clássico Master Mind, desenvolvido como projeto da disciplina de Programação para Web. O jogo foi construído utilizando tecnologias web modernas como React, Next.js e TypeScript, com uma interface responsiva inspirada na estética retrô 8-bit. O projeto aborda conceitos fundamentais de desenvolvimento web, incluindo gerenciamento de estado, componentização, responsividade e implementação de lógica de jogo. A aplicação oferece dois modos de jogo: singleplayer contra IA e multiplayer local, proporcionando uma experiência completa e interativa aos usuários. Os resultados demonstram a aplicabilidade das tecnologias web modernas na criação de jogos interativos com design responsivo e experiência de usuário aprimorada.
\end{abstract}

\begin{IEEEkeywords}
Master Mind, Desenvolvimento Web, React, Next.js, TypeScript, Jogos Digitais, Interface Responsiva, Estética 8-bit, Programação Web.
\end{IEEEkeywords}}

%%------------------------------------------------------
%% Início do Documento
%%------------------------------------------------------
\begin{document}

\maketitle
\IEEEdisplaynontitleabstractindextext
\IEEEpeerreviewmaketitle

%%------------------------------------------------------
%% Section: Introdução
%%------------------------------------------------------
\section{Introdução}

\IEEEPARstart{O}{Master Mind} é um jogo clássico de lógica e dedução criado em 1970 que desafia jogadores a descobrirem uma sequência secreta de cores através de tentativas sucessivas e feedback lógico. Este trabalho apresenta uma implementação digital moderna deste jogo clássico, desenvolvida como projeto prático da disciplina de Programação para Web na Universidade Federal de Santa Catarina (UFSC).

O objetivo principal deste projeto é aplicar conceitos fundamentais de desenvolvimento web front-end, utilizando tecnologias modernas e boas práticas de programação. A implementação foi realizada utilizando o framework React em conjunto com Next.js, TypeScript para tipagem estática, e CSS moderno para estilização, resultando em uma aplicação web completa, responsiva e interativa.

Uma característica distintiva deste projeto é a adoção de uma estética visual inspirada nos jogos retrô 8-bit dos anos 1980, criando uma experiência nostálgica que conecta o jogo clássico com sua versão digital moderna. Esta escolha estética não apenas homenageia a era de ouro dos videogames, mas também demonstra habilidades avançadas de design e implementação CSS.

O trabalho está organizado da seguinte forma: a Seção II apresenta os objetivos do projeto; a Seção III descreve os materiais e métodos utilizados; a Seção IV detalha a implementação técnica; a Seção V apresenta os resultados obtidos; e a Seção VI conclui o trabalho com considerações finais.

%%------------------------------------------------------
%% Section: Objetivos
%%------------------------------------------------------
\section{Objetivos}

\subsection{Objetivo Geral}
Desenvolver uma aplicação web interativa do jogo Master Mind utilizando tecnologias web modernas, demonstrando competências em programação front-end, design responsivo e implementação de lógica de jogo.

\subsection{Objetivos Específicos}
\begin{itemize}
\item Implementar o jogo Master Mind seguindo suas regras clássicas
\item Utilizar React e Next.js para construção da interface e gerenciamento de estado
\item Aplicar TypeScript para garantir tipagem estática e robustez do código
\item Desenvolver interface responsiva compatível com diferentes dispositivos
\item Criar sistema de IA para modo singleplayer
\item Implementar modo multiplayer local
\item Desenvolver estética visual retrô 8-bit consistente
\item Aplicar HTML5 semântico e acessibilidade
\item Implementar sistema de pontuação e feedback visual
\end{itemize}

%%------------------------------------------------------
%% Section: Materiais e Métodos
%%------------------------------------------------------
\section{Materiais e Métodos}

\subsection{Tecnologias Utilizadas}

A implementação do projeto utilizou as seguintes tecnologias:

\begin{itemize}
\item \textbf{Next.js 15.5.3}: Framework React para aplicações web com renderização híbrida
\item \textbf{React 19.1.0}: Biblioteca JavaScript para construção de interfaces de usuário
\item \textbf{TypeScript}: Superset JavaScript com tipagem estática
\item \textbf{Zustand 5.0.8}: Biblioteca para gerenciamento de estado global
\item \textbf{Tailwind CSS 4.1.13}: Framework CSS utility-first
\item \textbf{ESLint 9}: Ferramenta de análise estática de código
\item \textbf{Material Tailwind}: Componentes UI baseados em Material Design
\end{itemize}

\subsection{Arquitetura do Sistema}

A arquitetura do sistema segue o padrão de componentes React, organizada em três camadas principais:

\begin{enumerate}
\item \textbf{Camada de Apresentação}: Componentes React responsáveis pela interface do usuário
\item \textbf{Camada de Lógica}: Implementação das regras do jogo e algoritmos
\item \textbf{Camada de Estado}: Gerenciamento centralizado do estado da aplicação usando Zustand
\end{enumerate}

\subsection{Metodologia de Desenvolvimento}

O desenvolvimento seguiu uma abordagem iterativa e incremental:

\begin{enumerate}
\item \textbf{Fase 1 - Planejamento}: Definição de requisitos e design da arquitetura
\item \textbf{Fase 2 - Prototipação}: Desenvolvimento de protótipos da interface
\item \textbf{Fase 3 - Implementação}: Codificação dos componentes e lógica de jogo
\item \textbf{Fase 4 - Estilização}: Aplicação da estética 8-bit e responsividade
\item \textbf{Fase 5 - Testes}: Validação funcional e correção de bugs
\item \textbf{Fase 6 - Documentação}: Elaboração da documentação técnica
\end{enumerate}

%%------------------------------------------------------
%% Section: Implementação
%%------------------------------------------------------
\section{Implementação}

\subsection{Estrutura de Componentes}

A aplicação foi estruturada seguindo princípios de componentização e reutilização de código. Os principais componentes implementados incluem:

\subsubsection{Componentes da Página Inicial}
\begin{itemize}
\item \textbf{Header}: Cabeçalho com título e navegação
\item \textbf{Sidebar}: Barra lateral com informações do projeto
\item \textbf{Footer}: Rodapé com informações da disciplina
\item \textbf{GamePreview}: Prévia visual do jogo
\end{itemize}

\subsubsection{Componentes do Jogo}
\begin{itemize}
\item \textbf{SingleplayerGame}: Implementação do modo single player
\item \textbf{MultiplayerGame}: Implementação do modo multiplayer
\item \textbf{ModeSelector}: Seletor de modo de jogo
\end{itemize}

\subsection{Lógica do Jogo}

A implementação da lógica do Master Mind segue as regras clássicas:

\begin{enumerate}
\item Um jogador (ou IA) define uma sequência secreta de 4 cores
\item O outro jogador tenta adivinhar a sequência
\item Após cada tentativa, recebe feedback:
   \begin{itemize}
   \item Pinos pretos: cores corretas em posições corretas
   \item Pinos brancos: cores corretas em posições incorretas
   \end{itemize}
\item O jogo termina quando a sequência é descoberta ou após 10 tentativas
\end{enumerate}

\subsection{Sistema de IA}

Para o modo singleplayer, foi implementado um algoritmo de IA baseado em:

\begin{itemize}
\item Geração aleatória da sequência secreta
\item Validação de tentativas do jogador
\item Cálculo de feedback baseado em comparação de arrays
\item Sistema de pontuação baseado em número de tentativas
\end{itemize}

\subsection{Design Responsivo}

A interface foi desenvolvida com mobile-first approach, utilizando:

\begin{itemize}
\item Media queries CSS para diferentes breakpoints
\item Flexbox e Grid para layouts fluidos
\item Tipografia responsiva com unidades relativas
\item Componentes adaptáveis a diferentes tamanhos de tela
\end{itemize}

\subsection{Estética 8-bit}

A estética retrô foi implementada através de:

\begin{itemize}
\item Fonte 'Press Start 2P' do Google Fonts
\item Paleta de cores inspirada em consoles clássicos
\item Efeitos de glow com box-shadow e text-shadow
\item Animações pixeladas e scanlines
\item Bordas e elementos com estilo retrô
\end{itemize}

%%------------------------------------------------------
%% Section: Resultados
%%------------------------------------------------------
\section{Resultados}

\subsection{Funcionalidades Implementadas}

O projeto resultou em uma aplicação web totalmente funcional com as seguintes características:

\begin{itemize}
\item Interface responsiva funcionando em desktop, tablet e mobile
\item Dois modos de jogo completos (singleplayer e multiplayer)
\item Sistema de pontuação funcional
\item Feedback visual imediato para tentativas
\item Estética 8-bit consistente em toda aplicação
\item Navegação intuitiva entre seções
\item Performance otimizada com Next.js
\end{itemize}

\subsection{Página Inicial}

A página inicial implementa:

\begin{itemize}
\item Seção "Sobre" com descrição do projeto
\item Seção "Objetivos" listando metas do trabalho
\item Seção "Como Jogar" com instruções detalhadas
\item Seção "Regras" explicando mecânicas do jogo
\item Seção "Tecnologias" apresentando stack utilizada
\item Layout em três colunas (header, main, sidebar)
\item Links de navegação funcionais
\end{itemize}

\subsection{Experiência do Usuário}

A aplicação oferece uma experiência de usuário aprimorada através de:

\begin{itemize}
\item Transições suaves entre estados
\item Feedback visual claro para ações
\item Interface intuitiva e autoexplicativa
\item Responsividade em todos dispositivos
\item Estética nostálgica e envolvente
\end{itemize}

\subsection{Qualidade do Código}

O código fonte apresenta:

\begin{itemize}
\item Tipagem completa com TypeScript
\item Componentes reutilizáveis e modulares
\item Separação clara de responsabilidades
\item Código limpo e bem documentado
\item Padrões consistentes de nomenclatura
\item Configuração adequada de linting
\end{itemize}

%%------------------------------------------------------
%% Section: Conclusão
%%------------------------------------------------------
\section{Conclusão}

Este trabalho apresentou a implementação bem-sucedida de uma versão digital do jogo clássico Master Mind utilizando tecnologias web modernas. Os objetivos propostos foram alcançados, resultando em uma aplicação funcional, responsiva e esteticamente atraente.

A experiência de desenvolvimento proporcionou aprendizado significativo em:

\begin{itemize}
\item Desenvolvimento front-end com React e Next.js
\item Gerenciamento de estado em aplicações complexas
\item Implementação de lógica de jogo em JavaScript/TypeScript
\item Design responsivo e mobile-first
\item Estilização avançada com CSS moderno
\item Boas práticas de programação e arquitetura de software
\end{itemize}

A escolha da estética 8-bit demonstrou-se acertada, criando uma identidade visual única e uma experiência nostálgica que complementa o jogo clássico. O design responsivo garante acessibilidade em diversos dispositivos, expandindo o público potencial.

Como trabalhos futuros, sugere-se:

\begin{itemize}
\item Implementação de ranking online de pontuações
\item Modo multiplayer online com WebSockets
\item Níveis de dificuldade variáveis
\item Sistema de conquistas e troféus
\item Temas visuais alternativos
\item Suporte a internacionalização (i18n)
\end{itemize}

O projeto cumpriu seu papel educacional, proporcionando experiência prática com tecnologias relevantes do mercado e demonstrando a aplicabilidade dos conceitos de Programação para Web na criação de aplicações interativas e envolventes.

%%------------------------------------------------------
%% References
%%------------------------------------------------------
\bibliographystyle{IEEEtran}
\bibliography{references}

\end{document}